\documentclass[11pt]{article}

\usepackage{amsmath, amssymb, amsthm}
\usepackage{physics}
\usepackage{geometry}
\usepackage{hyperref}
\usepackage{mathtools}
\usepackage{enumitem}

\geometry{margin=1in}

% ------------------------------------------------------------------
% Theorem environments
% ------------------------------------------------------------------
\newtheorem{theorem}{Theorem}[section]
\newtheorem{lemma}[theorem]{Lemma}
\newtheorem{proposition}[theorem]{Proposition}
\newtheorem{corollary}[theorem]{Corollary}

\theoremstyle{definition}
\newtheorem{definition}[theorem]{Definition}
\newtheorem{example}[theorem]{Example}
\newtheorem{remark}[theorem]{Remark}

% ------------------------------------------------------------------
% Title
% ------------------------------------------------------------------
\title{%
	{\Large \textbf{Article 3}}\\[6pt]
	\textbf{Physical Interpretation of Twin Kernel Spaces:\\
		Relativistic, Quantum, and Geometric Perspectives on Kernel Transport}
}

\author{J.~Nembe}

\date{}

\begin{document}
	
	\maketitle
	
	\begin{abstract}
		This article develops the physical interpretation of Twin Kernel Spaces introduced in Articles~1 and~2.
		Kernel transport, represented mathematically as a group action on spaces of functions, is shown to correspond to physical notions of reference-frame change, Lorentz transformations, Hamiltonian flows, and quantum propagators.
		We establish that statistical smoothing can be interpreted as evolution under a geometric or physical flow, and that classical invariance principles in mathematical physics correspond to stability properties of kernel-based estimators.
		Connections with relativistic geometry, quantum evolution, gauge transformations, and renormalization flows are explored in a unified operator-theoretic framework.
	\end{abstract}
	
	\tableofcontents
	
	% ================================================================
	\section{Introduction}
	% ================================================================
	
	\subsection{Motivations from statistical geometry and physics}
	
	\subsection{Twin Kernel Spaces as a geometric theory of information transport}
	
	\subsection{Physical heuristics behind kernel transport: distance, time, and propagation}
	
	\subsection{Overview of contributions}
	
% ================================================================
\section{Twin kernels as geometric transformations}
\label{sec:twin-geometry}
% ================================================================

Twin Kernel Spaces arise when a base kernel is transported by a group action.
This section develops the geometric meaning of this transport.
We show that group actions behave as generalized reference-frame changes, that kernels transform like metrics under pullback, and that Twin Spaces form natural bundles of transported geometries.
This interpretation connects statistical kernel transport with the classical machinery of differential geometry and mathematical physics.

% ----------------------------------------------------------------
\subsection{Group actions as generalized reference-frame transformations}
% ----------------------------------------------------------------

Let $E$ be a smooth manifold and $G$ a Lie group acting smoothly on $E$.
For $g\in G$ we denote the action $x\mapsto g\cdot x$ and the tangent map $D(g)_x:T_xE\to T_{g\cdot x}E$.

\begin{definition}[Reference-frame transformation]
	A transformation $g\in G$ is interpreted as a change of reference frame, mapping each point $x\in E$ to $g\cdot x$ and pushing forward vectors, densities, and functions accordingly.
\end{definition}

In classical physics, such transformations include:
\begin{itemize}[itemsep=1pt,topsep=3pt]
	\item translations,
	\item rotations and Lorentz boosts,
	\item affine or conformal transformations,
	\item symplectic diffeomorphisms in Hamiltonian mechanics.
\end{itemize}

Twin Kernel Spaces elevate this physical intuition to an operator-theoretic framework.

\begin{definition}[Unitary transport]
	The action of $G$ induces a unitary operator
	\[
	(U_g f)(x)
	= J_{g}(x)^{-1/2}\, f(\overline{g}\cdot x),
	\]
	where $J_g$ is the Jacobian determinant.
\end{definition}

Thus $U_g$ implements the change of coordinate system associated with $g$.

% ----------------------------------------------------------------
\subsection{Transported kernels as metric deformations}
% ----------------------------------------------------------------

Let $K_e$ be a base reproducing kernel with RKHS $\mathcal{H}_e$.
For any $g\in G$, the transported kernel is defined by
\[
K_g(x,y)
= U_g K_e(\cdot,\cdot) U_g^{-1}(x,y)
= J_g(x)^{1/2} K_e(\overline{g}\cdot x,\overline{g}\cdot y)
J_g(y)^{1/2}.
\]

\begin{proposition}[Kernel transport as metric pullback]
	The kernel $K_g$ is the pullback of $K_e$ by the diffeomorphism $g$:
	\[
	K_g = (g^{-1})^\ast K_e.
	\]
	In particular, $K_g$ induces the pullback metric
	\[
	\langle u,v\rangle_{g,x}
	= \langle D(g^{-1})_x u,\, D(g^{-1})_x v\rangle_{e,g^{-1}\cdot x}.
	\]
\end{proposition}

\begin{proof}
	A direct computation using the pushforward–pullback identities for integral kernels and the invariance of $U_g$.
\end{proof}

Thus, the geometry underlying the RKHS is deformed by the group action.
Statistically, this corresponds to estimating in a coordinate system whose metric structure has been transported by $g$.

% ----------------------------------------------------------------
\subsection{Local diffeomorphisms, Lie groups, and infinitesimal generators}
% ----------------------------------------------------------------

When $G$ is a Lie group, the infinitesimal structure of kernel transport is governed by its Lie algebra $\mathfrak{g}$.

Let $g_t=\exp(tX)$ denote the flow of an element $X\in\mathfrak{g}$.
Then
\[
U_{g_t} = \exp(t \mathcal{L}_X),
\]
where $\mathcal{L}_X$ is the Lie derivative acting on functions.

\begin{proposition}[Infinitesimal kernel deformation]
	The derivative of the transported kernel at $t=0$ is
	\[
	\frac{d}{dt}K_{g_t}\Big|_{t=0}
	= \mathcal{L}_X K_e + (1/2)(\operatorname{div}X)\,K_e.
	\]
\end{proposition}

This identity mirrors the standard transformation of Riemannian metrics under flows:
\[
\frac{d}{dt}g_t = \mathcal{L}_X g.
\]

The analogy is exact: kernels play the role of generalized metrics.

% ----------------------------------------------------------------
\subsection{Twin spaces as bundles of transported metrics}
% ----------------------------------------------------------------

\begin{definition}[Twin space]
	The family of RKHSs
	\[
	\mathcal{H}_g := \mathcal{H}(K_g), \qquad g\in G,
	\]
	is called the \emph{twin space bundle}.
\end{definition}

The total space
\[
\mathcal{B} = \bigcup_{g\in G} \{g\}\times \mathcal{H}_g
\]
is a fiber bundle over $G$, with fiber $\mathcal{H}_e$ and transition maps given by $U_g$.

\begin{theorem}[Unitary trivialization of the twin bundle]
	The map
	\[
	G\times \mathcal{H}_e \to \mathcal{B},\qquad (g,f)\mapsto (g,U_g f)
	\]
	is a bundle isomorphism.
\end{theorem}

This result shows that Twin Kernel Spaces behave exactly like:
- the tangent bundle of a manifold,
- a gauge bundle in physics,
- a family of Hilbert spaces related by unitary parallel transport.

From this structural viewpoint, kernel smoothing corresponds to the application of a spectral operator within a fiber, while estimation across different geometries corresponds to transporting between fibers.

\begin{remark}
	Mathematically, this aligns Twin Kernel Spaces with the framework of fibrewise spectral geometry, used in quantum field theory on curved spacetimes.
\end{remark}

% ----------------------------------------------------------------
\subsection{Examples: affine, conformal, and Lorentz groups}
% ----------------------------------------------------------------

\begin{example}[Affine group $\mathrm{Aff}(d)$]
	For transformations $x\mapsto Ax + b$,
	the transported kernel becomes
	\[
	K_g(x,y) = |\det A|^{1/2}
	K_e(A^{-1}(x-b), A^{-1}(y-b))
	|\det A|^{1/2}.
	\]
	Gaussian kernels remain Gaussian; orthogonal polynomials transform by known affine recursions.
\end{example}

\begin{example}[Conformal group]
	If $g$ is conformal with factor $\Omega_g(x)$, then
	\[
	K_g(x,y)
	= \Omega_g(x)^{d/2}\, K_e(g^{-1}\cdot x, g^{-1}\cdot y) \,\Omega_g(y)^{d/2},
	\]
	mirroring the Weyl rescaling of the metric in conformal geometry.
\end{example}

\begin{example}[Lorentz group]
	For Minkowski spacetime with metric $\eta$ and Lorentz transformation $\Lambda$,
	kernel transport preserves the causal structure:
	\[
	K_\Lambda(x,y)
	= K_e(\Lambda^{-1}x,\Lambda^{-1}y).
	\]
	This will be central in Section~\ref{sec:relativity}.
\end{example}

These examples illustrate that kernel transport embeds a wide class of geometric and physical transformations in a unified operator-theoretic language.


	
	% ================================================================
	% ================================================================
	\section{Relativistic interpretation of kernel transport}
	\label{sec:relativity}
	% ================================================================
	
	Twin Kernel Spaces admit a natural relativistic interpretation.
	Under Lorentz transformations, the structure of the Minkowski metric is preserved, and transported kernels behave precisely as covariant two-point functions.
	This section formalizes the relationship between Lorentz invariance, hyperbolic geometry, and kernel transport, leading to a relativistic theory of statistical smoothing.
	
	% ----------------------------------------------------------------
	\subsection{Lorentz transformations as twin transports}
	% ----------------------------------------------------------------
	
	Consider Minkowski space $(\mathbb{R}^{1+d},\eta)$ with metric
	\[
	\eta = \mathrm{diag}(-1,1,\ldots,1).
	\]
	Let $\Lambda\in \mathrm{SO}(1,d)$ be a Lorentz transformation, i.e.
	\[
	\eta = \Lambda^\top \eta \Lambda.
	\]
	The induced action on spacetime is $x\mapsto \Lambda x$.
	
	\begin{definition}[Lorentz transport]
		For a base kernel $K_e$ on Minkowski space, its Lorentz-transported version is
		\[
		K_\Lambda(x,y) := K_e(\Lambda^{-1}x,\Lambda^{-1}y).
		\]
	\end{definition}
	
	This is precisely the twin kernel construction with $G=\mathrm{SO}(1,d)$.
	No Jacobian factor appears because Lorentz transformations preserve volume.
	
	\begin{proposition}[Lorentz invariance of transported kernels]
		If $K_e$ is Lorentz invariant, i.e.\ $K_e(\Lambda x,\Lambda y)=K_e(x,y)$ for all $\Lambda$, then
		\[
		K_\Lambda = K_e.
		\]
	\end{proposition}
	
	Lorentz invariance thus corresponds to fixed points of the twin transport under $G$.
	
	% ----------------------------------------------------------------
	\subsection{Twin kernels and Minkowski structures}
	% ----------------------------------------------------------------
	
	For arbitrary $K_e$, the transformed kernel $K_\Lambda$ behaves as a \emph{covariant two-point function}:
	\[
	K_\Lambda(x,y) 
	= K_e(\Lambda^{-1}x,\Lambda^{-1}y)
	= K_e(x',y')
	\]
	with coordinates $(x',y')$ measured in a different inertial frame.
	
	\begin{theorem}[Covariance of twin kernel metrics]
		Let $g_{K_e}$ denote the metric induced by an RKHS kernel $K_e$.
		Then the corresponding metric for the transported kernel satisfies
		\[
		g_{K_\Lambda}(x)
		= \Lambda^{-\top} g_{K_e}(\Lambda^{-1}x)\,\Lambda^{-1}.
		\]
	\end{theorem}
	
	\begin{proof}
		The induced metric is a covariant $2$-tensor, $g_{K_\Lambda}=(\Lambda^{-1})^\ast g_{K_e}$; with $K_\Lambda(x,y)=K_e(\Lambda^{-1}x,\Lambda^{-1}y)$ and the linear map $D(\Lambda^{-1})=\Lambda^{-1}$ this gives $g_{K_\Lambda}(x)=(\Lambda^{-1})^\top g_{K_e}(\Lambda^{-1}x)\,\Lambda^{-1}$. For instance, for $K_e(x,y)=\exp(-\tfrac12\|x-y\|^2)$ one has $g_{K_e}\equiv I$ and $g_{K_\Lambda}=\Lambda^{-\top}\Lambda^{-1}$, as required.
	\end{proof}

	Thus twin geometry reproduces the fundamental transformation law of tensor fields under Lorentz symmetry. (Note that the covariant law uses $\Lambda^{-\top}$ and $\Lambda^{-1}$; only for orthogonal transformations, where $\Lambda^{-1}=\Lambda^\top$, does it reduce to $\Lambda\,g\,\Lambda^\top$. For genuine boosts $\Lambda^{-1}\neq\Lambda^\top$, so the inverse transpose is essential.)
	
	\begin{remark}
		This makes the RKHS a natural “internal spacetime” whose geometry transforms covariantly under changes of inertial frame.
	\end{remark}
	
	% ----------------------------------------------------------------
	\subsection{Velocity composition, rapidity, and hyperbolic geometry}
	% ----------------------------------------------------------------
	
	Consider one-dimensional special relativity ($d=1$).
	Lorentz boosts are parametrized by the rapidity $\theta$, with matrix
	\[
	\Lambda(\theta)=
	\begin{pmatrix}
		\cosh\theta & \sinh\theta\\
		\sinh\theta & \cosh\theta
	\end{pmatrix}.
	\]
	Velocities compose by
	\[
	v = \tanh\theta.
	\]
	
	\begin{proposition}[Hyperbolic composition as twin transport]
		Let $g_\theta$ denote the twin transport induced by $\Lambda(\theta)$.
		Then composition of transports satisfies
		\[
		g_{\theta_1}\circ g_{\theta_2} = g_{\theta_1+\theta_2}.
		\]
		In particular,
		\[
		v_1 \oplus v_2
		= \tanh(\theta_1+\theta_2)
		= \frac{v_1+v_2}{1+v_1 v_2},
		\]
		the usual relativistic velocity addition.
	\end{proposition}
	
	\begin{proof}
		Follows from $\Lambda(\theta_1)\Lambda(\theta_2)=\Lambda(\theta_1+\theta_2)$ and the definition of $v=\tanh\theta$.
	\end{proof}
	
	Thus, the composition of observers corresponds precisely to the composition of twin transports.
	
	\begin{remark}
		Statistically, this means that changing coordinates twice corresponds to changing geometry by the composed hyperbolic parameter.
		This allows a unified analysis of anisotropic or accelerated coordinate systems.
	\end{remark}
	
	% ----------------------------------------------------------------
	\subsection{Relativistic invariants and statistical invariants}
	% ----------------------------------------------------------------
	
	In relativity, the spacetime interval
	\[
	s^2 = -t^2 + \|x\|^2
	\]
	is invariant under $\Lambda$.
	Analogously, in Twin Kernel Spaces the following are invariant:
	
	\begin{enumerate}[itemsep=3pt]
		\item the RKHS norm of a transported function (unitary invariance),
		\item the spectrum of the integral operator $T_e$ (unitary conjugation),
		\item spectral filters $\Gamma_k$,
		\item the induced pseudo-distance between probability measures.
	\end{enumerate}
	
	\begin{theorem}[Statistical invariance under Lorentz transport]
		For any $\Lambda\in \mathrm{SO}(1,d)$ and any probability measure $\mu$ on spacetime,
		\[
		\| T_\Gamma^{(\Lambda)} U_\Lambda f \|_{\mathcal{H}_\Lambda}
		= \| T_\Gamma f \|_{\mathcal{H}_e}.
		\]
	\end{theorem}
	
	\begin{proof}
		Using $T_\Gamma^{(\Lambda)}=U_\Lambda T_\Gamma U_\Lambda^{-1}$ and the unitarity of $U_\Lambda$.
	\end{proof}
	
	Thus smoothing is frame-independent when measured in the appropriate transported geometry.
	
	% ----------------------------------------------------------------
	\subsection{Twin estimation in expanding or collapsing geometries}
	% ----------------------------------------------------------------
	
	Relativistic transport naturally models changes of scale comparable to cosmological expansion.
	Let
	\[
	g_t=\Lambda(\theta(t))
	\]
	with $\theta(t)$ increasing.
	Then the transported kernel evolves as
	\[
	K_{g_t}(x,y) = K_e(\Lambda(-\theta(t)) x,\Lambda(-\theta(t)) y).
	\]
	
	\begin{example}[Expanding reference frame]
		If $\theta(t)\sim Ht$ (constant expansion rate),
		then the effective geometry “stretches” at exponential rate.
		The corresponding estimator
		\[
		\widehat f_t = T_\Gamma^{(g_t)} \widehat f_{\mathrm{emp},t}
		\]
		adapts automatically to the expanding geometry.
	\end{example}
	
	\begin{example}[Light-cone smoothing]
		If $K_e$ has support within a fixed light-cone, its transported version also respects causality.
		This yields statistical procedures respecting the Minkowski structure.
	\end{example}
	
	Such connections suggest that Twin Kernel Spaces naturally model:
	- accelerated observers,
	- expanding/contracting geometric regimes,
	- relativistic signal processing,
	- causal estimation.
	
	
	
% ================================================================
\section{Quantum interpretation and propagators}
\label{sec:quantum}
% ================================================================

Quantum mechanics is fundamentally a theory of linear operators acting on Hilbert spaces.
Twin Kernel Spaces offer a natural operator-theoretic setting in which the evolution of quantum states, the propagation of wave functions, and the structure of Green’s functions appear as special cases of kernel transport.
This section develops the quantum interpretation of transported kernels as propagators and shows that statistical smoothing corresponds to imaginary-time evolution.

% ----------------------------------------------------------------
\subsection{Unitary group actions as quantum evolutions}
% ----------------------------------------------------------------

In quantum mechanics, the state of a system lives in a Hilbert space $\mathcal{H}$ and evolves according to a one-parameter unitary group
\[
U_t = e^{-i t H},
\]
generated by the Hamiltonian $H$.

In Twin Kernel Spaces, the group $G$ acting on the domain $E$ induces a unitary representation
\[
U_g : \mathcal{H}_e \to \mathcal{H}_g,
\qquad (U_g f)(x) = J_g(x)^{-1/2} f(g^{-1}\cdot x),
\]
analogous to a change of quantum reference frame.

\begin{proposition}[Twin transports as generalized Schrödinger evolutions]
	Let $g_t=\exp(tX)$ be the flow of $X\in\mathfrak{g}$.
	Then
	\[
	U_{g_t} = e^{t \mathcal{L}_X},
	\]
	where $\mathcal{L}_X$ is the Lie derivative.
	Thus $U_{g_t}$ plays the role of a quantum evolution operator with generator $\mathcal{L}_X$.
\end{proposition}

This identifies $\mathcal{L}_X$ as an analogue of a quantum Hamiltonian.

% ----------------------------------------------------------------
\subsection{Spectral decomposition and twin eigenmodes}
% ----------------------------------------------------------------

Let $K_e$ be a Mercer kernel with eigen-decomposition
\[
K_e(x,y)=\sum_{k\ge 0}\lambda_k \phi_k(x)\phi_k(y).
\]
As shown in Articles~1–2,
transported kernels satisfy
\[
K_g(x,y) = \sum_{k\ge 0}\lambda_k\,\phi_k^{(g)}(x)\phi_k^{(g)}(y),
\qquad
\phi_k^{(g)} = U_g\phi_k.
\]

\begin{proposition}[Twin eigenmodes as quantum stationary states]
	The functions $\phi_k^{(g)}$ form an orthonormal basis of $\mathcal{H}_g$, and obey
	\[
	T_e^{(g)} \phi_k^{(g)} = \lambda_k \phi_k^{(g)},
	\qquad T_e^{(g)} := U_g T_e U_g^{-1}.
	\]
	(In general $T_e\phi_k^{(g)}\neq\lambda_k\phi_k^{(g)}$ unless $U_g$ commutes with $T_e$, i.e.\ unless $g$ is a symmetry of the base kernel.)
	Thus the twin eigenfunctions play the role of stationary states transported into a different quantum frame.
\end{proposition}

The analogy is exact: eigenfunctions are energy eigenstates, and transport changes the representation of the same physical mode.

% ----------------------------------------------------------------
\subsection{Kernel smoothing as imaginary-time propagation}
% ----------------------------------------------------------------

In quantum mechanics, imaginary-time evolution
\[
e^{-tH}
\]
corresponds to \emph{Wick rotation}, replacing $t\mapsto it$.
It transforms unitary evolution into diffusion and suppresses high-frequency modes.

Kernel smoothing in twin spaces achieves exactly the same effect.

Let
\[
T_\Gamma = \sum_{k} \Gamma_k\, \phi_k\otimes \phi_k,
\quad
\Gamma_k\in[0,1].
\]

Choosing $\Gamma_k = e^{-t\lambda_k}$ yields a heat semigroup.

\begin{theorem}[Twin smoothing = imaginary-time quantum evolution]
	For any transported geometry $g$,
	\[
	T_\Gamma^{(g)} = U_g e^{-t T_e} U_g^{-1}.
	\]
	Thus smoothing by $T_\Gamma^{(g)}$ is imaginary-time propagation under the “Hamiltonian” $T_e$ expressed in the reference frame $g$.
\end{theorem}

\begin{proof}
	Using $T_\Gamma^{(g)} = U_g T_\Gamma U_g^{-1}$ and the spectral decomposition of $T_\Gamma$.
\end{proof}

This provides a physical explanation for the effectiveness of spectral filtering.

% ----------------------------------------------------------------
\subsection{Twin kernels and Green's functions}
% ----------------------------------------------------------------

In quantum field theory, a Green’s function $G(x,y)$ satisfies
\[
(H+m^2) G(\cdot,y) = \delta_y.
\]

Similarly, the resolvent of the transported operator $T_e^{(g)} = U_g T_e U_g^{-1}$
plays the role of a Green's function: for $\alpha>0$,
\[
(T_e^{(g)} + \alpha I)\, G_\alpha^{(g)}(\cdot,y) = \delta_y,
\]
which yields the spectral (resolvent) representation below.

\begin{proposition}[Twin kernels as Green’s functions]
	For any $\alpha>0$, the kernel
	\[
	G_\alpha^{(g)}(x,y) = \sum_k (\lambda_k+\alpha)^{-1}\phi_k^{(g)}(x)\phi_k^{(g)}(y)
	\]
	is the Green’s function of the operator $T_e+\alpha I$ transported by $g$.
\end{proposition}

Thus twin transport preserves the physical role of kernels as propagators.

% ----------------------------------------------------------------
\subsection{Gauge transformations and twin coordinate changes}
% ----------------------------------------------------------------

In gauge theory, a field $\psi$ transforms by
\[
\psi \mapsto e^{i\chi} \psi,
\]
and potentials transform by gradients of $\chi$.

Twin coordinate changes behave analogously.

If $g$ acts on $E$, and $U_g$ acts on functions by precomposition with $g^{-1}$, then
\[
f^{(g)}(x)=U_g f(x)
\]
is the analogue of a gauge-transformed field.

\begin{proposition}[Gauge covariance]
	For any spectral operator $T_\Gamma$,
	\[
	T_\Gamma^{(g)} = U_g T_\Gamma U_g^{-1}
	\]
	is the gauge-transformed version of $T_\Gamma$.
\end{proposition}

This is the mathematical expression of gauge covariance in Twin Kernel Spaces.

% ----------------------------------------------------------------
\subsection{Summary of the quantum interpretation}
% ----------------------------------------------------------------

The quantum interpretation can be summarized as follows:
\begin{itemize}[itemsep=2pt]
	\item Transport $U_g$ corresponds to a quantum change of representation.
	\item Eigenfunctions $\phi_k^{(g)}$ are stationary states transported by $g$.
	\item Smoothing $T_\Gamma^{(g)}$ corresponds to imaginary-time evolution.
	\item Green’s functions arise from spectral resolvents.
	\item Gauge covariance emerges naturally from the twin transport rule.
\end{itemize}

This formalizes the deep structural analogy between kernel-based statistical estimation, quantum evolution, and gauge-covariant operator theory.


	
% ================================================================
\section{Hamiltonian flows and geometric transport}
\label{sec:hamiltonian}
% ================================================================

Hamiltonian dynamics provides a canonical geometric mechanism for transporting functions and probability densities.
In this section we show that Twin Kernel Spaces naturally encode Hamiltonian flows: kernel transport corresponds to the pullback of observables under a canonical transformation, while smoothing operators correspond to averaging along Hamiltonian trajectories.
This establishes a deep connection between statistical transport, classical mechanics, and the geometric theory of information.

% ----------------------------------------------------------------
\subsection{Hamiltonian dynamics as kernel transport}
% ----------------------------------------------------------------

Let $(M,\omega)$ be a symplectic manifold with Hamiltonian function $H:M\to\mathbb{R}$.
The Hamiltonian vector field $X_H$ is defined by
\[
\iota_{X_H}\omega = dH.
\]
Its flow $(\Phi_t)_{t\in\mathbb{R}}$ satisfies
\[
\frac{d}{dt}\Phi_t(x) = X_H(\Phi_t(x)).
\]

The flow consists of canonical transformations:
\[
\Phi_t^\ast\omega = \omega.
\]

\begin{definition}[Hamiltonian kernel transport]
	Given a base kernel $K_e$, we define
	\[
	K_t(x,y):=K_e(\Phi_{-t}(x),\Phi_{-t}(y)).
	\]
\end{definition}

This is exactly the twin kernel construction for the group $G=\{\Phi_t: t\in\mathbb{R}\}$.

\begin{proposition}[Symplectic invariance]
	If $K_e$ depends only on symplectic invariants (e.g.\ radial Gaussians, action-integrals), then $K_t=K_e$ for all $t$.
\end{proposition}

Thus Hamiltonian invariants correspond to statistical invariants under Hamiltonian twin transport.

% ----------------------------------------------------------------
\subsection{Twin kernels as generating functions of canonical transformations}
% ----------------------------------------------------------------

The key geometric observation is that any canonical transformation can be described by a generating function.
Twin kernel transport reproduces this mechanism.

Let $S_t:M\times M\to \mathbb{R}$ be a generating function of $\Phi_t$.
Then
\[
p = \partial_x S_t(x,y),\qquad
q = \partial_y S_t(x,y)
\]
encodes $\Phi_t:(x,p)\mapsto (y,q)$.

\begin{theorem}[Kernels encode generating functions]
	Assume $K_e(x,y)=\kappa(S_0(x,y))$ for some analytic $\kappa$.
	Then
	\[
	K_t(x,y)=\kappa(S_t(x,y)).
	\]
	Thus the twin-transported kernel propagates the generating function along the Hamiltonian flow.
\end{theorem}

\begin{proof}
	Follows from the identity $S_t(x,y)=S_0(\Phi_{-t}(x),\Phi_{-t}(y))$.
\end{proof}

Hence kernels act as “action-like” quantities transported along classical trajectories.

% ----------------------------------------------------------------
\subsection{Symplectic invariance and transport in RKHS geometry}
% ----------------------------------------------------------------

Hamiltonian flows preserve volumes:
\[
\det D\Phi_t(x)=1,
\]
and are therefore represented by unitary operators in $L^2(M)$:
\[
(U_t f)(x)=f(\Phi_{-t}(x)).
\]

\begin{proposition}[Unitarity of Hamiltonian transport in RKHS]
	The operator $U_t$ is unitary on $\mathcal{H}_e$ if and only if the kernel is symplectically invariant:
	\[
	K_e(\Phi_t x,\Phi_t y)= K_e(x,y)
	\quad\forall t.
	\]
\end{proposition}

Transport in $\mathcal{H}_e$ along $\Phi_t$ thus plays the role of a classical phase-space evolution.

\begin{remark}
	This mirrors the Koopman-von Neumann formulation of classical mechanics, where Hamiltonian evolution acts as a unitary operator on an $L^2$-space of classical observables.
	Twin Kernel Spaces provide a generalization adapted to nonlinear geometry and non-Euclidean metrics.
\end{remark}

% ----------------------------------------------------------------
\subsection{Energy, action, and information transport}
% ----------------------------------------------------------------

Hamiltonian flows conserve energy:
\[
H\circ \Phi_t = H.
\]
Thus the “energy landscape” of the estimation problem remains invariant under twin transport.

\begin{theorem}[Action invariance under twin transport]
	Let $\gamma$ be a trajectory of $\Phi_t$ and let
	\[
	\mathcal{A}(\gamma)=\int_\gamma p\,dq - H\,dt
	\]
	be its classical action.
	If $K_e$ is an action-type kernel,
	\[
	K_e(x,y) = F(\mathcal{A}(\gamma_{x,y})),
	\]
	then
	\[
	K_t(x,y)=K_e(x,y).
	\]
\end{theorem}

\begin{proof}
	Since $\Phi_t$ preserves $H$ and the 1-form $p\,dq$, the action integral is invariant.
\end{proof}

Thus kernels associated with action integrals define statistically invariant geometries under Hamiltonian transport.

\begin{remark}[Information transport]
	Under twin transport, statistical distances evolve as
	\[
	d_{K_t}(f,g)=d_{K_e}(f\circ\Phi_t, g\circ \Phi_t).
	\]
	Thus information is transported along Hamiltonian trajectories without distortion.
\end{remark}

% ----------------------------------------------------------------
\subsection{Examples and physical interpretations}
% ----------------------------------------------------------------

\begin{example}[Quadratic Hamiltonians: Gaussian kernels]
	Let $H(q,p)=\frac12(p^2+q^2)$.
	Then $\Phi_t$ is a rotation:
	\[
	(q,p)\mapsto (q\cos t + p\sin t\, ,\, p\cos t - q\sin t).
	\]
	A Gaussian kernel
	\[
	K_e(x,y)=\exp(-\|x-y\|^2/2\sigma^2)
	\]
	is invariant:
	\[
	K_t(x,y)=K_e(x,y).
	\]
	This mirrors the fact that harmonic motion preserves Euclidean distances in phase space.
\end{example}

\begin{example}[Integrable systems]
	For separable Hamiltonians $H(p,q)=H_1(p)+H_2(q)$, the flow decomposes and twin kernels factorize.
	The transported geometry corresponds to shearing the coordinate system along invariant tori.
\end{example}

\begin{example}[Chaotic flows]
	For Hamiltonians generating Anosov flows, twin transport produces exponentially diverging directions, inducing strong anisotropy in transported kernels.
	This provides a statistical analogue of stable/unstable manifolds.
\end{example}

% ----------------------------------------------------------------
\subsection{Summary}
% ----------------------------------------------------------------

Hamiltonian flows provide a physical and geometric model of twin kernel transport:
\begin{itemize}[itemsep=3pt]
	\item Transport corresponds to canonical pullback by $\Phi_t$.
	\item Eigenfunctions behave like stationary states transported along trajectories.
	\item Generating functions propagate through the kernel.
	\item Symplectic invariance corresponds to fixed points of transport.
	\item Information flows rigidly along Hamiltonian paths.
\end{itemize}

This bridges classical mechanics, spectral operator theory, and kernel-based inference into a unified geometric formalism.

	
% ================================================================
\section{Diffusion, renormalization, and scale interactions}
\label{sec:diffusion-renorm}
% ================================================================

Twin Kernel Spaces naturally encode nonlinear and multiscale geometric transformations.
This section shows that diffusion semigroups, renormalization flows, and heat kernel regularization appear as special cases of spectral filtering transported by a group action.
We develop a unifying viewpoint in which smoothing, coarse-graining, and scale interactions are governed by transported Laplacians and their spectral properties.

% ----------------------------------------------------------------
\subsection{Heat flow as smoothing flow}
% ----------------------------------------------------------------

Let $(M,g)$ be a Riemannian manifold with Laplace–Beltrami operator $\Delta_g$.
The heat semigroup
\[
P_t = e^{t\Delta_g}
\]
satisfies
\[
\partial_t u_t = \Delta_g u_t,
\qquad
u_0=f.
\]

In Twin Kernel Spaces, the analogue of $P_t$ is obtained by taking the spectral filter
\[
\Gamma_k(t) = e^{-t\lambda_k},
\]
which yields the operator
\[
T_t f = \sum_{k\ge 0} e^{-t\lambda_k}\,\langle f,\phi_k\rangle\,\phi_k.
\]

\begin{proposition}[Twin heat flow]
	For any transported geometry $g\in G$,
	\[
	T_t^{(g)} = U_g T_t U_g^{-1}
	\]
	is the heat flow in the twin geometry $(M,g\cdot g)$.
\end{proposition}

Thus the heat equation in one geometry immediately induces heat evolution in all twin geometries.

\begin{remark}
	This establishes a direct link between smoothing in statistics and diffusion in geometric analysis.
\end{remark}

% ----------------------------------------------------------------
\subsection{Diffusion semigroups and transported Laplacians}
% ----------------------------------------------------------------

Let $g\in G$ be a smooth group action on $M$.
Define the transported Laplacian
\[
\Delta_g := U_g \Delta_e U_g^{-1}.
\]

\begin{theorem}[Twin diffusion semigroup]
	The family
	\[
	P_t^{(g)} := e^{t\Delta_g}
	\]
	forms a diffusion semigroup on the transported geometry:
	\[
	\partial_t u_t = \Delta_g u_t.
	\]
	Moreover,
	\[
	P_t^{(g)} = U_g P_t U_g^{-1}.
	\]
\end{theorem}

\begin{proof}
	Direct spectral calculus: since $\Delta_g$ and $\Delta_e$ are unitarily conjugate, their semigroups are conjugate as well.
\end{proof}

This shows that diffusion is invariant under twin transport: the heat equation retains its structure under any geometry deformation induced by $G$.

\begin{remark}
	In geometric analysis, $U_g$ plays the role of a pullback operator for heat kernels.
	In Twin Kernel Spaces, it acts as a geometric change of “observer”.
\end{remark}

% ----------------------------------------------------------------
\subsection{Renormalization flows and progressive smoothing}
% ----------------------------------------------------------------

Renormalization group (RG) theory describes progressive coarse-graining of a system by integrating out small scales.
In the spectral domain, coarse-graining corresponds to suppressing high eigenvalues.

Let
\[
T_\ell = \sum_k \Gamma_k(\ell)\,\phi_k\otimes \phi_k,
\qquad
\Gamma_k(\ell) = \exp(-\ell\,\lambda_k),
\]
where $\ell$ is the RG scale.

Then
\[
T_{\ell_1+\ell_2} = T_{\ell_1}T_{\ell_2},
\]
corresponding to a classical RG flow law.

\begin{theorem}[Twin renormalization flow]
	Define the renormalization operator in the geometry transported by $g$:
	\[
	T_\ell^{(g)} = U_g T_\ell U_g^{-1}.
	\]
	Then $\{T_\ell^{(g)}\}_{\ell\ge 0}$ forms a renormalization semigroup in the transported geometry.
\end{theorem}

\begin{proof}
	Unitary conjugation preserves the semigroup property.
\end{proof}

\begin{remark}
	This formalizes the idea that renormalization is a spectral coarse-graining procedure invariant under changes of observer or geometry.
\end{remark}

% ----------------------------------------------------------------
\subsection{Multiscale twin geometries and information propagation}
% ----------------------------------------------------------------

Twin Kernel Spaces naturally encode multiscale geometry, since the generator of scale evolution is fully spectral.

Let $(\lambda_k)$ be eigenvalues of $T_e$ ordered increasingly.
Small scales correspond to large $k$.

Consider a scale parameter $s$ and define
\[
g_s = \exp(sX),
\]
where $X\in\mathfrak{g}$ generates a geometric deformation.
Then the transported geometry $g_s\cdot E$ exhibits scale-dependent metric properties.

\begin{proposition}[Scale–geometry interaction]
	The effective kernel at scale $(\ell,s)$ is
	\[
	K_{\ell,s}(x,y)
	= \sum_{k} e^{-(\ell\lambda_k)} \phi_k^{(g_s)}(x)\phi_k^{(g_s)}(y).
	\]
\end{proposition}

Thus two effects combine:
\begin{itemize}
	\item spectral decay $e^{-\ell\lambda_k}$ (coarse-graining),
	\item geometric deformation $\phi_k^{(g_s)}=U_{g_s}\phi_k$.
\end{itemize}

\begin{theorem}[Information propagation along twin flows]
	Let $f$ be a signal transported by the geometric flow $g_s$ and smoothed by scale $\ell$.
	Then
	\[
	T_\ell^{(g_s)} f = U_{g_s} T_\ell (f\circ g_s^{-1}).
	\]
	Information propagation thus separates into:
	\[
	\text{geometry transport} \quad+\quad \text{scale transport}.
	\]
\end{theorem}

This result reveals a fundamental decoupling: geometry affects eigenfunctions, while renormalization affects eigenvalues.

% ----------------------------------------------------------------
\subsection{Examples: heat kernels, Wilson RG, and conformal flows}
% ----------------------------------------------------------------

\begin{example}[Heat kernel flows]
	On a manifold with curvature $R$, the heat kernel
	\[
	K_t(x,y) \sim \frac{1}{(4\pi t)^{d/2}}
	\exp\!\left(-\frac{d_g(x,y)^2}{4t}\right)(1+\mathcal{O}(t))
	\]
	is transported by $g$ via the pullback of geodesic distances.
	Twin transport acts as a “deformed heat kernel”.
\end{example}

\begin{example}[Wilson renormalization]
	Let $\ell$ denote the Wilson scale.
	Then the flow of effective kernels satisfies
	\[
	K_\ell = e^{-\ell T_e} K_0,
	\]
	and transported versions satisfy
	\[
	K_\ell^{(g)} = U_g e^{-\ell T_e} U_g^{-1}K_0.
	\]
	This corresponds to rewriting the RG flow in a deformed geometry.
\end{example}

\begin{example}[Conformal heat flows]
	If $g$ is a conformal transformation with factor $\Omega_g$,
	then the transported Laplacian satisfies
	\[
	\Delta_g = \Omega_g^{-2}\Delta_e + \text{lower-order terms}.
	\]
	The resulting heat flow corresponds to a conformal RG flow.
\end{example}

% ----------------------------------------------------------------
\subsection{Summary}
% ----------------------------------------------------------------

Diffusion and renormalization processes fit naturally into Twin Kernel Spaces:
\begin{itemize}[itemsep=3pt]
	\item Diffusion is smoothing via spectral decay.
	\item Renormalization is coarse-graining via suppression of large eigenvalues.
	\item Both processes are invariant under twin geometry transport.
	\item Geometry acts on eigenfunctions via unitaries $U_g$.
	\item Scale acts on eigenvalues via $(\Gamma_k)$.
	\item The combined geometry–scale flow describes how information propagates across scales.
\end{itemize}

This formalizes the intuition that smoothing, diffusion, and renormalization are deeply linked geometric processes governed by spectral transport.

	
% ================================================================
\section{Twin geometries in curved spaces and field theories}
\label{sec:curved}
% ================================================================

In this section we extend the framework of Twin Kernel Spaces to curved manifolds and classical field theories.
We show that kernel transport behaves as a pullback of the underlying geometric structures:
metrics, Laplace–Beltrami operators, propagators, and stress–energy tensors.
This leads to a geometric–operator formulation of twin estimation on general spacetimes.

% ----------------------------------------------------------------
\subsection{Twin transport on Riemannian manifolds}
% ----------------------------------------------------------------

Let $(M,g)$ be a smooth Riemannian manifold and $K_e$ a Mercer kernel compatible with $g$.
For a diffeomorphism $\Phi:M\to M$, define the transported kernel
\[
K_\Phi(x,y) := K_e(\Phi^{-1}(x),\Phi^{-1}(y)).
\]

\begin{proposition}[Pullback interpretation]
	The kernel $K_\Phi$ is the pullback of $K_e$ by $\Phi$:
	\[
	K_\Phi = (\Phi^{-1})^\ast K_e.
	\]
\end{proposition}

Thus twin transport generalizes the covariant action of diffeomorphisms on tensor fields.

Let $\nabla$ be the Levi–Civita connection.
The Laplace–Beltrami operator is
\[
\Delta_g f = \mathrm{div}_g(\nabla f).
\]

\begin{definition}[Transported Laplacian]
	The twin geometry associated with $\Phi$ induces the operator
	\[
	\Delta_\Phi := U_\Phi \,\Delta_g\, U_\Phi^{-1},
	\]
	where $U_\Phi f = f\circ\Phi^{-1}$.
\end{definition}

This operator acts as the Laplacian in the transported geometry $(M,\Phi^\ast g)$.

\begin{theorem}[Twin Laplacian covariance]
	For all smooth $f$,
	\[
	\Delta_\Phi (f \circ \Phi) = (\Delta_g f)\circ \Phi.
	\]
\end{theorem}

\begin{proof}
	This is the standard coordinate-covariance of the Laplace–Beltrami operator, expressed via a unitary conjugation.
\end{proof}

Thus twin kernels capture the essential covariant structure of Riemannian geometry.

% ----------------------------------------------------------------
\subsection{Curvature, geodesics, and kernel deformation}
% ----------------------------------------------------------------

On a curved manifold, the heat kernel satisfies the short-time expansion
\[
K_t(x,y)
= \frac{\exp\!\left(-\frac{d_g(x,y)^2}{4t}\right)}{(4\pi t)^{d/2}}
\left[ 1 + \frac{t}{6} R(x) + \mathcal{O}(t^2)\right],
\]
where $R$ is the scalar curvature.

Transporting the geometry by $\Phi$ produces
\[
d_{\Phi^\ast g}(x,y)=d_g(\Phi^{-1}x,\Phi^{-1}y),
\qquad
R_{\Phi^\ast g}(x)=R(g)(\Phi^{-1}x).
\]

\begin{proposition}[Twin curvature transport]
	Heat-based kernels satisfy
	\[
	K_t^{(\Phi)}(x,y)
	=K_t(\Phi^{-1}x,\Phi^{-1}y)
	\]
	with identical curvature corrections, expressed in the transported coordinates.
\end{proposition}

Thus curvature information is transported coherently through twin kernels, preserving the geometric meaning.

\begin{remark}
	Geodesic deviation under $\Phi$ becomes the geometric analogue of anisotropic kernel deformation.
\end{remark}

% ----------------------------------------------------------------
\subsection{Field equations and kernel propagators}
% ----------------------------------------------------------------

Classical fields on $(M,g)$ satisfy equations of the form
\[
\mathcal{D}_g \psi = J,
\]
where:
- $\mathcal{D}_g$ is a differential operator (Klein–Gordon, Dirac, Maxwell),
- $J$ is a source,
- $\psi$ is a field.

Green’s functions satisfy
\[
\mathcal{D}_g G_g(\cdot,y) = \delta_y.
\]

\begin{definition}[Twin propagator]
	Given a diffeomorphism $\Phi$, define
	\[
	G_\Phi(x,y) := G_g(\Phi^{-1}x,\Phi^{-1}y).
	\]
\end{definition}

\begin{theorem}[Twin covariance of field propagators]
	For any differential operator $\mathcal{D}_g$,
	\[
	\mathcal{D}_{\Phi^\ast g}\, G_\Phi(\cdot,y)
	= \delta_y.
	\]
\end{theorem}

\begin{proof}
	Apply unitary covariance:  
	$\mathcal{D}_{\Phi^\ast g} = U_\Phi \mathcal{D}_g U_\Phi^{-1}$ and 
	$G_\Phi = U_\Phi G_g U_\Phi^{-1}$.
\end{proof}

Hence twin propagators correspond exactly to field propagators in transformed geometries.

% ----------------------------------------------------------------
\subsection{Twin kernels and stress-energy interpretations}
% ----------------------------------------------------------------

The stress–energy tensor of a field $\psi$ is
\[
T_{\mu\nu} = 
\partial_\mu\psi\,\partial_\nu\psi
- \frac12 g_{\mu\nu}\big( g^{\alpha\beta}\partial_\alpha\psi\partial_\beta\psi + m^2 \psi^2 \big).
\]

For kernels, the analogue is the “energy density”
\[
\mathcal{E}_{K_g}(x) 
= \sum_{k} \lambda_k\,|\phi_k^{(g)}(x)|^2,
\]
interpreted as the contribution of spectral modes to local variance.

\begin{proposition}[Twin stress-energy transport]
	Under twin transport,
	\[
	\mathcal{E}_{K_{\Phi}}(x)
	= \mathcal{E}_{K_e}(\Phi^{-1}x).
	\]
\end{proposition}

\begin{proof}
	Since $\phi_k^{(\Phi)} = U_\Phi \phi_k$, the result follows from unitarity.
\end{proof}

Thus the analogues of local field-energy densities transform covariantly under twin geometry.

\begin{remark}
	This interpretation links local statistical complexity (variance of modes) to the stress–energy distribution of a physical field.
\end{remark}

% ----------------------------------------------------------------
\subsection{Examples: de Sitter, Schwarzschild, and wave kernels}
% ----------------------------------------------------------------

\begin{example}[de Sitter spacetime]
	Let $(M,g)$ be de Sitter space with scale factor $a(t)$.
	Twin transport by a diffeomorphism rescaling $t$ yields a transported geometry with modified expansion.
	Kernel propagators correspond to mode functions in different cosmological gauges.
\end{example}

\begin{example}[Schwarzschild geometry]
	In Schwarzschild spacetime, the radial coordinate transformation
	\[
	r\mapsto r^\prime = r + 2m\log(r-2m)
	\]
	induces a twin transport $\Phi$ that straightens radial null geodesics.
	Twin kernels correspond to propagators in Eddington–Finkelstein coordinates.
\end{example}

\begin{example}[Wave kernels on curved spaces]
	For the wave operator $\square_g$, the Hadamard parametrix
	\[
	G(x,y) \sim \Delta(x,y)^{1/2}\,\delta(\sigma(x,y))
	\]
	transports under $\Phi$ by transforming both Synge’s world function $\sigma$ and the Van Vleck determinant $\Delta$.
\end{example}

% ----------------------------------------------------------------
\subsection{Summary}
% ----------------------------------------------------------------

Twin geometries on curved spaces unify geometric and field-theoretic transformations:
\begin{itemize}[itemsep=3pt]
	\item Kernel transport behaves as pullback of metrics and Laplacians.
	\item Curvature and geodesic structure deform kernels covariantly.
	\item Field propagators (Green’s functions) transform under twin transport.
	\item Stress–energy analogues in RKHSs follow identical covariance laws.
	\item The framework applies to arbitrary spacetimes (de Sitter, Schwarzschild, FLRW).
\end{itemize}

This places the Twin Kernel framework in direct alignment with the mathematics of geometric field theories.

	
% ================================================================
\section{Applications to physics-inspired statistical models}
\label{sec:applications}
% ================================================================

The geometric and physical interpretation of Twin Kernel Spaces enables the construction of new classes of statistical models that inherit structural invariances from physics.
In this section we present several applications: relativistic regression models, quantum-inspired kernel machines, twin flows in nonlinear dynamical systems, and transport-based regularization in learning.

These examples demonstrate the practical implications of the theoretical formalism developed in Sections~2--7.

% ----------------------------------------------------------------
\subsection{Relativistic regression and invariant densities}
% ----------------------------------------------------------------

In special relativity, physical laws must be covariant under Lorentz transformations.
Twin kernel transport provides a natural mechanism to construct statistical estimators that satisfy the same invariance principles.

Let $(x_i,y_i)$ be observed data and let $g\in G=\mathrm{SO}(1,d)$ act via Lorentz transformations.
Define the transported regression estimator
\[
\widehat f_g = T_\Gamma^{(g)} \widehat f_{\mathrm{emp}}.
\]

\begin{definition}[Lorentz-invariant regression model]
	A regression model is Lorentz invariant if
	\[
	\widehat f_g(\Lambda x) = \widehat f_e(x)
	\quad\text{for all }\Lambda\in\mathrm{SO}(1,d).
	\]
\end{definition}

\begin{theorem}[Covariance of twin regression]
	If the data transform as
	\[
	(x_i,y_i)\mapsto (\Lambda x_i, y_i),
	\]
	and the base kernel is Lorentz covariant, then
	\[
	\widehat f_\Lambda = U_\Lambda\widehat f_e.
	\]
\end{theorem}

\begin{proof}
	By the conjugation identity $T_\Gamma^{(\Lambda)} = U_\Lambda T_\Gamma U_\Lambda^{-1}$ and the fact that $U_\Lambda$ acts by precomposition.
\end{proof}

Thus, relativistic regression is a direct application of twin geometry.

\begin{remark}
	This suggests new statistical procedures for high-energy data, particle trajectories, cosmic-ray distributions, and time-like/spacelike correlated signals.
\end{remark}

% ----------------------------------------------------------------
\subsection{Quantum-inspired kernel machines}
% ----------------------------------------------------------------

Quantum mechanics suggests that data may be represented as superpositions of eigenmodes, with amplitudes evolving in time.
Twin Kernel Spaces recover this structure naturally.

Let
\[
f(x)=\sum_k a_k \phi_k(x)
\]
be the expansion of a signal in the transported eigenbasis.

\begin{definition}[Quantum-inspired kernel machine]
	Let $U(g)$ be a unitary representation of $G$.
	The quantum-inspired estimator is
	\[
	\widehat f_{\mathrm{quant}}(x)
	= \sum_k \Gamma_k\,a_k\, \phi_k^{(g)}(x),
	\]
	where $\Gamma_k$ acts as a spectral gate (analogous to a projector or propagator).
\end{definition}

If $\Gamma_k(t)=e^{-it\lambda_k}$, then the estimator evolves according to
\[
\widehat f_t = e^{-it T_e}\widehat f_0,
\]
interpretable as a Schrödinger evolution in RKHS.

\begin{proposition}[Imaginary/real time duality]
	Replacing $\Gamma_k(t)=e^{-it\lambda_k}$ by $\Gamma_k(t)=e^{-t\lambda_k}$ yields a transition from quantum evolution to diffusion smoothing, corresponding to Wick rotation.
\end{proposition}

This duality suggests novel learning algorithms that interpolate between unitary and dissipative regimes.

% ----------------------------------------------------------------
\subsection{Twin flows in dynamical systems}
% ----------------------------------------------------------------

Twin transports generalize the Koopman operator framework for deterministic dynamical systems.

Let $(M,\Phi_t)$ be a flow on a manifold and $f\in L^2(M)$ an observable.
The Koopman operator is
\[
U_t f = f\circ\Phi_{-t}.
\]

In Twin Kernel Spaces:
\[
f^{(t)} = U_t f,\qquad
K_t(x,y)=K_e(\Phi_{-t}(x),\Phi_{-t}(y)).
\]

\begin{theorem}[Twin-Koopman equivalence]
	For any smooth flow,
	\[
	T_\Gamma^{(t)} f
	= U_t T_\Gamma\,(f\circ\Phi_t),
	\]
	and the twin-spectral estimator evolves exactly as an observable under Koopman dynamics.
\end{theorem}

\begin{proof}
	Conjugation by $U_t$ combined with the spectral calculus for $T_\Gamma$.
\end{proof}

This provides a natural representation of nonlinear dynamical systems as geometric deformations of kernel spaces.

\begin{example}[Chaotic dynamics]
	For Anosov flows, transported kernels exhibit anisotropies aligned with stable/unstable manifolds.
	This mirrors the hyperbolic structure of the underlying dynamics.
\end{example}

% ----------------------------------------------------------------
\subsection{Transport-based regularization}
% ----------------------------------------------------------------

Regularization can be interpreted as penalizing energy, curvature, or oscillation.
Twin kernel transport provides new regularizers of the form
\[
\mathcal{R}(f)
= \|T_\Gamma^{(g)} f\|_{\mathcal{H}_g}^2,
\]
which combine scale filtering and geometric transformation.

\begin{proposition}[Geometry-sensitive regularization]
	If $g$ encodes a deformation of the domain (stretching, shear, contraction), then $\mathcal{R}(f)$ penalizes oscillations measured in the transported metric $(g^\ast g)$.
\end{proposition}

Such regularizers are sensitive to:
\begin{itemize}[itemsep=2pt]
	\item anisotropies in physical systems,
	\item directional flows in time series,
	\item non-Euclidean or curved coordinate systems,
	\item areas of high curvature or compression.
\end{itemize}

\begin{example}[Physical priors]
	If $g$ encodes a relativistic contraction/expansion, $\mathcal{R}(f)$ incorporates Lorentz-type priors.
	If $g_t$ corresponds to a Hamiltonian flow, $\mathcal{R}(f)$ encodes invariance under canonical transformations.
\end{example}

% ----------------------------------------------------------------
\subsection{Applications to temporal processes and signal transport}
% ----------------------------------------------------------------

Twin flows provide a powerful language for time-varying processes.

Let $g_t$ be a time-dependent deformation (velocity field, linear acceleration, dilation).
Define the transported estimator
\[
\widehat f_t = T_\Gamma^{(g_t)} \widehat f_{\mathrm{emp},t}.
\]

\begin{proposition}[Adaptive smoothing in moving frames]
	If $g_t$ encodes a moving coordinate frame, then $\widehat f_t$ performs smoothing in that frame, equivalent to a physically natural “motion-compensated” estimator.
\end{proposition}

This applies to:
\begin{itemize}[itemsep=2pt]
	\item video streams and deforming grids,
	\item sound propagation (Doppler effects),
	\item nonlinear time-warped signals,
	\item data observed by accelerating or rotating sensors.
\end{itemize}

\begin{example}[Doppler-aware estimation]
	Let $g_t$ be a Lorentz-type dilation encoding Doppler shift.
	Then $\widehat f_t$ tracks the signal in the correct frequency–time geometry without requiring explicit model fitting.
\end{example}

% ----------------------------------------------------------------
\subsection{Summary}
% ----------------------------------------------------------------

Twin Kernel Spaces provide a broad and unified family of statistical models inspired by physics:
\begin{itemize}[itemsep=3pt]
	\item Lorentz-invariant regression and density estimation.
	\item Quantum-inspired learning machines with unitary or dissipative dynamics.
	\item Koopman-type representations of nonlinear dynamical systems.
	\item Geometry-based regularization reflecting physical deformations.
	\item Adaptive smoothing in moving, accelerating, or curved frames.
\end{itemize}

These applications illustrate the versatility of the twin framework and its ability to incorporate deep structural invariances from physical theories into statistical learning.

	
% ================================================================
\section{Conclusion and future directions}
\label{sec:conclusion}
% ================================================================

This article has developed a physical and geometric interpretation of Twin Kernel Spaces.
Building on the foundations established in Articles~1 and~2, we have shown that kernel transport behaves as a generalized reference-frame change, incorporating the transformation laws of classical geometry, relativistic kinematics, quantum evolution, Hamiltonian mechanics, diffusion processes, and renormalization flows.
Twin Kernel Spaces thus unify notions from mathematical physics with spectral methods in nonparametric inference.

\subsection*{Summary of contributions}

The main contributions of this work may be summarized as follows:

\begin{itemize}[itemsep=3pt]
	\item \textbf{Geometric transport.}  
	Kernel transport is the pullback of tensor-like structures under group actions, producing families of RKHSs related by unitary conjugation. This provides a geometric bundle structure over the transformation group.
	
	\item \textbf{Relativistic interpretation.}  
	Lorentz transformations act as twin transports, reproducing hyperbolic composition laws, causal invariances, and Minkowski covariance. Statistical estimation acquires physical consistency across inertial frames.
	
	\item \textbf{Quantum interpretation.}  
	Twin transport generalizes Schrödinger evolution, with smoothing corresponding to imaginary-time propagation and Green’s functions realized as kernel resolvents in transported geometries.
	
	\item \textbf{Hamiltonian dynamics.}  
	Canonical transformations appear as kernel transports generated by Hamiltonian flows. Eigenfunctions behave as transported stationary states and generating functions propagate through the kernel geometry.
	
	\item \textbf{Diffusion, renormalization, and scales.}  
	Heat kernels, diffusion semigroups, and renormalization group flows arise from spectral filters transported by geometric actions. Scale and geometry decouple cleanly through eigenvalue–eigenfunction separation.
	
	\item \textbf{Field theories and curved spaces.}  
	On Riemannian or Lorentzian manifolds, twin kernels reproduce the transformation laws of propagators, Laplacians, curvature terms, and stress–energy analogues.
	
	\item \textbf{Physics-inspired statistical models.}  
	Relativistic regression, quantum-inspired kernel machines, geometric regularization, and adaptive smoothing in moving frames are obtained directly from the twin formalism.
\end{itemize}

Altogether, these results reveal Twin Kernel Spaces as a natural operator-theoretic language unifying statistical inference, geometric analysis, and mathematical physics.

\subsection*{Conceptual interpretation}

Twin Kernel Spaces suggest that statistical information does not live in a fixed geometry.
Instead, geometry itself depends on transformations of the observer, the coordinate system, or the physical regime.
The twin viewpoint systematically accounts for these deformations by transporting both metrics and eigenmodes.
In this light:

\begin{itemize}[itemsep=2pt]
	\item statistical estimation resembles physical propagation;
	\item smoothing resembles dissipative or Euclidean-time evolution;
	\item invariance corresponds to symmetries of the underlying geometry;
	\item regularization corresponds to coarse-graining or renormalization in scale.
\end{itemize}

This perspective opens new pathways for designing estimators by enforcing physical principles such as covariance, causality, or conservation laws.

\subsection*{Open problems}

The theory developed here raises several open questions:

\begin{enumerate}[itemsep=3pt]
	\item \textbf{Twin curvature estimation.}  
	Can curvature invariants be estimated consistently using twin-deformed Laplacians?
	
	\item \textbf{Nonlinear twin transports.}  
	How do transported geometries interact with nonlinear PDEs, e.g.\ conservation laws or reaction–diffusion systems?
	
	\item \textbf{Gauge fields and twin connections.}  
	Is there a natural analogue of a gauge connection on the twin bundle, enabling parallel transport between twin fibres?
	
	\item \textbf{Stochastic twin flows.}  
	Can stochastic differential equations be formulated in twin geometries, leading to new diffusion models?
	
	\item \textbf{Renormalization limits.}  
	Does long-time twin smoothing converge to universal fixed points in analogy with RG universality classes?
	
	\item \textbf{Statistical optimality.}
	Under which conditions do twin estimators achieve minimax or oracle properties across transported geometries? \emph{(Partially answered for $G$-invariant density estimation on a quotient $E/G$: the exact minimax rate $n^{-2s/(2s+d_{\mathrm{eff}})}$, $d_{\mathrm{eff}}=\dim(E/G)$, is attained by a projection estimator---which pays the model's metric complexity $D/n$ rather than a description-length coding cost $D\log n$---and is carried across transported geometries by the unitary invariance of the loss.)}
\end{enumerate}

Addressing these questions would deepen the connection between physical laws and statistical inference.

\subsection*{Directions for further research}

The twin framework suggests several promising research directions:

\begin{itemize}[itemsep=3pt]
	\item \textbf{Twin quantum machine learning.}  
	Designing learning architectures based on unitary dynamics, propagators, and spectral gates.
	
	\item \textbf{Twin geometry for dynamical data.}  
	Using time-varying transports $g_t$ to analyze signals, trajectories, and PDE-driven datasets.
	
	\item \textbf{Twin field theories.}  
	Developing kernel-based analogues of Klein–Gordon, Dirac, or Yang–Mills fields.
	
	\item \textbf{Twin statistical physics.}  
	Studying phase transitions in the spectral distribution of transported operators.
	
	\item \textbf{Applications to cosmology and spacetime data.}  
	Constructing Lorentz-covariant estimators for cosmological time series, particle events, or gravitational signals.
	
	\item \textbf{Twin inference on manifolds.}  
	Applying the formalism to data on curved spaces: shapes, biological surfaces, networks, and Lie groups.
\end{itemize}

\subsection*{Final perspective}

Twin Kernel Spaces offer a unified view of geometry, physics, and statistical inference.
They show that transporting kernels by group actions yields a rich family of geometric models whose physical analogues include inertial transformations, quantum evolutions, Hamiltonian flows, diffusion, and renormalization.
By connecting these structures, the twin framework opens a path toward a comprehensive geometric theory of inference—one in which estimation is inseparable from the geometry of information and the symmetries that govern it.

The development of such a theory promises new insights at the intersection of geometry, physics, and data science, and suggests that the fundamental principles of mathematical physics may provide a natural foundation for next-generation statistical methodologies.

	
	\appendix
\appendix

% ================================================================
\section{Technical Appendices}
\label{sec:appendix-technical}
% ================================================================

This appendix provides the operator identities, spectral calculi, geometric transformations, and detailed proofs that underpin the results of Sections~2–7.

% ================================================================
\subsection{Appendix A: Operator identities for twin transport}
% ================================================================

Let $G$ be a Lie group acting smoothly on a manifold $E$.
For $g\in G$, define the unitary operator
\[
(U_gf)(x)=J_g(x)^{-1/2} f(g^{-1}\cdot x).
\]

\begin{lemma}[Group property]
	\[
	U_{gh} = U_g U_h,\qquad U_e=I,\qquad U_g^{-1}=U_{g^{-1}}.
	\]
\end{lemma}

\begin{proof}
	Direct computation: $U_g U_h f(x) = J_g(x)^{-1/2} J_h(g^{-1}\cdot x)^{-1/2} f((gh)^{-1}\cdot x)$, so the representation property $U_{gh}=U_gU_h$ holds with the (half-density) cocycle $J_{gh}(x)=J_g(x)\,J_h(g^{-1}\cdot x)$, the chain rule for the Jacobian of $x\mapsto g^{-1}\cdot x$.
\end{proof}

\begin{lemma}[Adjoint]
	$U_g$ is unitary in $L^2(E)$ and therefore
	\[
	U_g^\ast = U_{g^{-1}}.
	\]
\end{lemma}

\begin{proof}
	Using the change of variables $y=g^{-1}\cdot x$ and the definition of $J_g$.
\end{proof}

\begin{proposition}[Transport of kernels]
	Let $K_e$ be a positive definite kernel.
	Define
	\[
	K_g(x,y) = J_g(x)^{1/2} K_e(g^{-1}x,g^{-1}y) J_g(y)^{1/2}.
	\]
	Then
	\[
	K_g = U_g K_e U_g^{-1}.
	\]
\end{proposition}

\begin{proof}
	Apply $U_g$ to $K_e(\cdot,y)$ and evaluate the reproducing property.
\end{proof}

% ================================================================
\subsection{Appendix B: Spectral calculus under group actions}
% ================================================================

Let
\[
K_e(x,y)=\sum_{k\ge 0}\lambda_k \phi_k(x)\phi_k(y)
\]
be a Mercer decomposition.

Transported eigenfunctions are
\[
\phi_k^{(g)} = U_g\phi_k.
\]

\begin{lemma}[Spectral transport]
	\[
	K_g(x,y)=\sum_{k\ge 0}\lambda_k\,\phi_k^{(g)}(x)\phi_k^{(g)}(y).
	\]
\end{lemma}

\begin{proof}
	Insert the decomposition into the definition of $K_g$ and use unitarity of $U_g$.
\end{proof}

\begin{proposition}[Transport of spectral operators]
	For any spectral filter $\Gamma=(\Gamma_k)$,
	\[
	T_\Gamma^{(g)} = U_g T_\Gamma U_g^{-1}.
	\]
\end{proposition}

\begin{proof}
	By direct action on $\phi_k^{(g)}$.
\end{proof}

This establishes the spectral covariance central to the twin framework.

% ================================================================
\subsection{Appendix C: Geometric transformations (pullback and pushforward)}
% ================================================================

Let $\Phi:E\to E$ be a diffeomorphism.  
For a function $f$, the pullback is
\[
\Phi^\ast f = f\circ \Phi.
\]

For vector fields:
\[
\Phi_\ast X = D\Phi\cdot X.
\]

\begin{proposition}[Pullback of metrics]
	If $g$ is a Riemannian metric, then
	\[
	(\Phi^\ast g)_x(u,v)
	= g_{\Phi(x)}(D\Phi_x u, D\Phi_x v).
	\]
\end{proposition}

\begin{lemma}[Transport of Laplacians]
	Let $\Delta_g$ be the Laplace–Beltrami operator. Then
	\[
	U_\Phi\, \Delta_g\, U_\Phi^{-1} = \Delta_{\Phi^\ast g}.
	\]
\end{lemma}

\begin{proof}
	Use the identity $\Phi^\ast(\mathrm{div}_g X) = \mathrm{div}_{\Phi^\ast g}(\Phi^\ast X)$ and the invariance of $\nabla$ under pullback.
\end{proof}

This matches the definition of the transported Laplacian $\Delta_\Phi$ in Section 7.

% ================================================================
\subsection{Appendix D: Variational formulas and infinitesimal generators}
% ================================================================

Let $g_t=\exp(tX)\in G$ with generator $X\in\mathfrak{g}$.

\begin{lemma}[Derivative of transported functions]
	\[
	\frac{d}{dt} U_{g_t} f\Big|_{t=0}= \mathcal{L}_X f,
	\]
	where $\mathcal{L}_X$ is the Lie derivative.
\end{lemma}

\begin{proof}
	Differentiate $f(g_t^{-1}x)$ at $t=0$.
\end{proof}

\begin{theorem}[Infinitesimal kernel deformation]
	\[
	\frac{d}{dt}K_{g_t}\big|_{t=0}
	= \mathcal{L}_X K_e + \frac12 (\mathrm{div}\,X)\,K_e.
	\]
\end{theorem}

\begin{proof}
	Differentiate the transported kernel definition and use the infinitesimal form of $J_{g_t}$.
\end{proof}

This formula is the kernel analogue of metric variation under a flow.

% ================================================================
\subsection{Appendix E: Relativistic covariance calculations}
% ================================================================

Let $\eta$ be the Minkowski metric and $\Lambda\in\mathrm{SO}(1,d)$ a Lorentz map.

\begin{lemma}[Lorentz metric invariance]
	\[
	\Lambda^\top \eta \Lambda = \eta.
	\]
\end{lemma}

\begin{proposition}[Covariance of twin metrics]
	Let $g_{K_e}$ be the RKHS metric induced by $K_e$.  
	Then
	\[
	g_{K_\Lambda}(x) = \Lambda^{-\top} g_{K_e}(\Lambda^{-1}x)\,\Lambda^{-1}.
	\]
\end{proposition}

\begin{proof}
	Direct computation using the reproducing property under transformed coordinates.
\end{proof}

\begin{lemma}[Velocity composition]
	If $\Lambda(\theta)$ is a boost with rapidity $\theta$,
	\[
	\Lambda(\theta_1)\Lambda(\theta_2) = \Lambda(\theta_1+\theta_2).
	\]
\end{lemma}

This gives the relativistic velocity addition law of Section 3.

% ================================================================
\subsection{Appendix F: Quantum evolution and Green’s functions}
% ================================================================

Let $H$ be a positive operator.  
Spectral calculus gives
\[
e^{-tH} = \sum_k e^{-t\lambda_k} \phi_k\otimes \phi_k.
\]

\begin{proposition}[Imaginary-time evolution]
	Twin smoothing with $\Gamma_k(t)=e^{-t\lambda_k}$ satisfies
	\[
	T_t^{(g)} = U_g e^{-tH} U_g^{-1}.
	\]
\end{proposition}

\begin{proof}
	Apply $U_g$ to each eigencomponent.
\end{proof}

For Green’s functions:
\[
G_\alpha = (H+\alpha I)^{-1} = \sum_k (\lambda_k+\alpha)^{-1}\phi_k\otimes\phi_k.
\]

\begin{lemma}[Twin propagator identity]
	\[
	G_\alpha^{(g)} = U_g G_\alpha U_g^{-1}.
	\]
\end{lemma}

This justifies the covariance of field propagators in Section 4.

% ================================================================
\subsection{Appendix G: Hamiltonian flows and symplectic geometry}
% ================================================================

Let $(M,\omega)$ be symplectic and let $\Phi_t$ be the Hamiltonian flow.

\begin{lemma}[Symplectic invariance]
	\[
	\Phi_t^\ast \omega = \omega.
	\]
\end{lemma}

\begin{lemma}[Volume preservation]
	\[
	\det D\Phi_t = 1.
	\]
\end{lemma}

\begin{proposition}[Hamiltonian unitarity]
	$U_t f = f\circ\Phi_{-t}$ is unitary.
\end{proposition}

\begin{proof}
	Use Liouville’s theorem for Hamiltonian flows.
\end{proof}

\begin{proposition}[Generating functions transported by kernels]
	If $K_e(x,y)=\kappa(S_0(x,y))$, then
	\[
	K_t(x,y)=\kappa(S_t(x,y)),
	\]
	where $S_t$ generates $\Phi_t$.
\end{proposition}

% ================================================================
\subsection{Appendix H: Diffusion, renormalization, and scale calculus}
% ================================================================

Let $\Delta$ be a Laplacian.  
Then the heat kernel satisfies
\[
K_t = e^{t\Delta}.
\]

Transport under $U_g$ gives:

\begin{lemma}[Transport of the heat semigroup]
	\[
	P_t^{(g)} = U_g P_t U_g^{-1}.
	\]
\end{lemma}

\begin{proposition}[Renormalization semigroup]
	If $T_\ell=e^{-\ell T_e}$, then
	\[
	T_\ell^{(g)} = U_g T_\ell U_g^{-1}
	\]
	and
	\[
	T_{\ell_1+\ell_2}^{(g)} = T_{\ell_1}^{(g)} T_{\ell_2}^{(g)}.
	\]
\end{proposition}

This matches the Wilson RG structure.

% ================================================================
\subsection{Appendix I: Curved spaces and field theory computations}
% ================================================================

Let $\Phi$ be a diffeomorphism.

\begin{lemma}[Curvature transport]
	\[
	R_{\Phi^\ast g}(x) = R_g(\Phi^{-1}x).
	\]
\end{lemma}

\begin{lemma}[Green’s functions under pullback]
	\[
	G_{\Phi}(x,y)=G_g(\Phi^{-1}x,\Phi^{-1}y).
	\]
\end{lemma}

\begin{proposition}[Stress–energy transport]
	If
	\[
	\mathcal{E}_g(x) = \sum_k \lambda_k|\phi_k(x)|^2,
	\]
	then
	\[
	\mathcal{E}_{g,\Phi}(x) = \mathcal{E}_g(\Phi^{-1}x).
	\]
\end{proposition}

\begin{proof}
	Use $\phi_k^{(\Phi)}=U_\Phi\phi_k$ and unitarity.
\end{proof}

These identities justify the results of Section~7 on field-theoretic covariance.


	
	\begin{thebibliography}{9}
	\bibitem{NembeTKS1} J.~Nembe, \emph{Twin Kernel Spaces I: Foundations} (Article~1).
	\bibitem{NembeTKS2} J.~Nembe, \emph{Twin Kernel Spaces II: Spectral Transport} (Article~2).
	\end{thebibliography}

\end{document}

